\section{扑克：Inference-Time Search 为什么可能压过训练规模}

Poker 是不完全信息博弈：玩家看不到对手手牌，行动还会反过来暴露信息。它因此特别适合检验一个问题：如果训练得到的策略已经很强，是否还值得在每个新局面上临时求解。讲者从 2015 年的人机赛失败讲起，不是为了列历史，而是为了说明研究文化曾经低估了推理时计算的收益。

\subsection{从 2015 年失利到“当场多想一会儿”}

第一张图把人类牌手与当时 AI 的差距放在同一张记分牌上。2015 年 Brains vs. AI 比赛中，Claudico 面对四位顶尖 heads-up no-limit Texas Hold'em 选手并未取胜。关键诊断不是“参数还不够多”，而是 bot 在困难局面仍近似查表式地瞬间行动；人类则会把更多时间分配给高风险手牌。

\lecturefigure{slide-01-2015-brains-ai-poker.jpg}{2015 Brains vs. AI Poker Competition：搜索转向前的失败基线}{官方录像恢复幻灯片 V001；Stanford CS25 Lecture 14}

\teachervoice{老师强调，最值得模仿的人类行为不是某套扑克技巧，而是\textbf{按局面难度动态分配思考时间}。容易的牌快速行动，决定整场收益的牌则值得重新求解。这个观察把研究问题从“再训练一个更大的策略”改写为“给现有策略多少 inference compute”。}

\subsection{用 exploitability 读懂搜索曲线}

扑克策略不能只看它对固定测试对手的胜率，因为一个策略可能恰好克制该对手，却能被另一种针对性打法轻易利用。更稳健的量是 \term{exploitability}：最优对手相对于均衡收益最多还能从该策略榨取多少额外收益。若用 $u_i(\pi_i,\pi_{-i})$ 表示玩家 $i$ 在联合策略下的期望收益，可把近似均衡误差写成
\begin{equation}
\epsilon_i(\pi)
= \max_{\pi_i'} u_i(\pi_i',\pi_{-i}) - u_i(\pi_i,\pi_{-i}).
\end{equation}
其中 $\pi_i'$ 是玩家 $i$ 的任意最佳响应，$\pi_{-i}$ 是其他玩家策略；$\epsilon_i$ 越小，说明策略越难被针对。课堂图纵轴写作 distance from Nash，本质上要读成“离不可利用策略还有多远”，而不是普通 accuracy。

\lecturefigure{slide-02-search-poker-scaling.jpg}{中等规模扑克中的 scaling：search 曲线整体显著低于 no-search}{官方录像恢复幻灯片 V002；曲线引用 Brown 与 Sandholm 的安全子博弈求解工作}

\begin{knowledgebox}{读图：先比较纵向差距，再看横向 scaling}
横轴 buckets 是 base abstraction / policy 的规模，从约 200 增至 20,000；纵轴是对数刻度的 Nash distance，越低越好。蓝线表示不做 search，橙线表示在当前子博弈中重新求解。两条线都会随 base model 变大而下降，但橙线在每个规模上都低很多。讲者把代表点概括为 exploitability 约下降七倍：真正强的信号不是某个点，而是 search 带来的整条曲线位移。
\end{knowledgebox}

\begin{warningbox}{曲线没有证明什么}
这张图不能证明“search 永远等价于把模型扩大固定倍数”，更不能把扑克中的七倍改进直接外推到语言任务。它只证明在该不完全信息博弈和该求解器下，局部推理时计算比继续扩张同类 base abstraction 更高效。
\end{warningbox}

\subsection{为什么这么有效的方法仍被忽视}

第三张 slide 不是技术细节，而是研究激励的诊断。许多团队把资源投入可复用的离线策略；search 需要每一步额外计算，工程路径复杂，且早期社区曾认为已知方法在大型游戏中不够有效。于是大家反复扩大 base model，却很少测量“同一模型在决策时多算一会儿”能换来多少收益。

\lecturefigure{slide-03-why-search-overlooked.jpg}{为什么 poker search 长期没有成为默认路线}{官方录像恢复幻灯片 V003；Stanford CS25 Lecture 14}

\teachervoice{讲者把这段历史讲成一个研究管理教训：当社区普遍相信某个方向“不会扩展”时，团队会停止为它建基础设施，缺少基础设施又进一步让实验看起来昂贵。search 的迟到既有算法原因，也有文化、算力预算与发表激励原因。}

\subsection{Libratus：把局部求解带进高水平单挑}

2017 年 Libratus 在 120,000 手牌的人机赛中战胜四位顶尖 heads-up no-limit poker 选手。这里的贡献不是每回合暴力枚举完整游戏树，而是先有一个离线 blueprint strategy，再在实际到达的子博弈里安全地细化策略。这样，系统把昂贵计算集中在真实出现、且对收益重要的分支上。

\lecturefigure{slide-04-libratus-2017.jpg}{Libratus：搜索驱动的双人无限注扑克系统}{官方录像恢复幻灯片 V004；Brown 与 Sandholm，Science 2017}

\begin{importantbox}{Blueprint 与局部 search 的分工}
Blueprint 提供全局可行的先验策略，使系统不会从零开始；subgame solving 则在已观察到的公共局面附近重新分配概率质量。前者负责覆盖，后者负责精度。若只有 blueprint，困难局面过于粗糙；若每次都从零求全局均衡，计算又不可承受。
\end{importantbox}

\subsection{Pluribus：从两人零和扩展到六人桌}

多人 poker 不再满足简洁的两人零和结构：某个行动对第三方的影响、临时利益一致和非传递性都会让求解更难。2019 年 Pluribus 仍以 modest compute 击败多位人类职业牌手，说明 search 的价值并不只存在于最干净的理论设定。图中最该读的不是杂志封面，而是“六人、实时、有限预算”这三个约束同时成立。

\lecturefigure{slide-05-pluribus-2019.jpg}{Pluribus：六人无限注德州扑克中的搜索系统}{官方录像恢复幻灯片 V005；Brown 与 Sandholm，Science 2019}

\begin{knowledgebox}{从扑克抽象出的通用 search 接口}
一个实用推理时求解器至少需要四件事：候选行动生成器、对未来状态的展开规则、状态或策略价值估计器、以及明确的预算停止条件。不同领域会替换具体算法，但“用 base model 提议，用 evaluator 比较，用预算控制深度”的接口可以迁移。
\end{knowledgebox}

\begin{lstlisting}[caption={Budgeted inference-time search 的领域无关骨架},label={lst:budgeted-search}]
def search(state, base_policy, transition, evaluator, budget):
    frontier = [(state, 1.0, [])]
    best = None
    for _ in range(budget):
        node = select_promising(frontier)
        for action, prior in base_policy(node.state):
            next_state = transition(node.state, action)
            value = evaluator(next_state)
            candidate = (value, node.path + [action])
            best = max(best, candidate) if best else candidate
            frontier.append(expand(node, next_state, prior, value))
    return best[1][0]
\end{lstlisting}

这段伪代码故意不指定 tree search、counterfactual regret minimization 或 beam search。它表达的是资源账户：预算增加时，系统可以探索更多候选、看得更深或调用更精确的 evaluator；如果 evaluator 有系统偏差，更多 search 也可能把偏差放大，因此计算量从来不是单独的质量保证。

\subsection{本章小结}

Poker 的核心证据是：同一 base policy 配上局部 search，可以显著降低可利用性；Libratus 和 Pluribus 则证明这种设计能在真实时间和有限硬件中工作。下一步要问的是，这条规律是否只属于扑克，还是能够上升为更一般的 inference-time computation 原则。

